%=======================================================================
\section{Results and Discussion}

The following points structure the interpretation of the benchmark and statistical validation
results.

\begin{itemize}[leftmargin=1.4em,itemsep=3pt]
  \item \textbf{Best overall configuration.} State which (feature-set $\times$ model) pair won on test
        macro-F1 and by how much, and confirm whether the full fusion beat the individual branches ---
        quantify the gain.
  \item \textbf{Branch contribution.} Compare A vs B vs C vs D and the cumulative sets. Did HOG
        dominate (it is the largest branch)? Did adding GLCM/LBP texture or DWT measurably help over
        HOG alone?
  \item \textbf{Accuracy vs.\ cost.} Reference the training-time-vs-F1 scatter: if a lighter model
        (e.g.\ SVM or LightGBM) matches the best within its confidence interval, argue for it on
        efficiency grounds.
  \item \textbf{Statistical significance.} Tie the leaderboard to Section~\ref{sec:stats}: are the top
        models statistically separable (McNemar / Friedman+Nemenyi), or is the top cluster a
        statistical tie? This is the most important point a grader looks for.
  \item \textbf{Per-class behaviour.} From the confusion matrices, note which tumor types are confused
        most (commonly Glioma $\leftrightarrow$ Meningioma) and relate it back to the texture
        similarity discussed in Section~1.
\end{itemize}

\subsection{Threats to validity}
\begin{itemize}[leftmargin=1.4em,itemsep=2pt]
  \item Single dataset and a single train/test split (mitigated by stratification, fixed seed, and CV
        --- supplemented by a seed-stability check across three random seeds).
  \item Phase~1 selects parameters on the training set using separability metrics, not downstream
        accuracy; the metrics are a fast proxy, so the chosen settings are good but not guaranteed
        globally optimal for every classifier.
  \item Handcrafted features cannot adapt the way a CNN would; the trade-off is interpretability,
        speed, and low data requirements.
\end{itemize}
